\section{Final Conclusion}

A small ODE model with \textbf{an LNP depot, cell-associated pre-cytosolic mRNA, cytosolic mRNA, and antigen protein} is sufficient to reproduce the key antigen-expression phenotypes observed in mRNA--LNP vaccination: delayed rise, a peak, and a decaying tail, with strong dependence on delivery and stability parameters.

The most interpretable mechanistic insight is that \textbf{total antigen exposure (AUC)} factorizes into a product of step efficiencies and lifetimes:
\begin{equation}
\mathrm{AUC}_P \propto D_0 \times f_{\mathrm{up}} \times f_{\mathrm{esc}} \times \frac{k_{\mathrm{tl}}}{k_{\mathrm{deg}}k_{\mathrm{clr}}},
\end{equation}
so small efficiencies \textemdash{} especially endosomal escape, widely reported to be low \textemdash{} can dominate platform performance.

Differences between sharp/transient vs.~lower/sustained expression at fixed dose can be explained by:
\begin{itemize}
    \item time-shaping of supply (burst vs.~sustained depot/uptake),
    \item mRNA and protein lifetimes, and
    \item innate-response--driven translation suppression, which can be included as a minimal inhibitory state when abrupt post-peak drops are observed.
\end{itemize}
